\rok{Фебруар 2026.}{А}

Други поправни тест из Алгоритама и структура података

Први практични тест одржан 19.02.2026.

\zadatak{1}{Конверзија листе у матрицу суседства}
Написати алгоритам који ће вршити директну конверзију листе суседства у матрицу суседства. Алгоритам као улазни параметар треба да прихвати динамички низ чворова, а треба да резултује креираном матрицом.

\textbf{Додатне напомене за решавање задатка:}
\begin{itemize}
	\item У template-у је закоментарисана метода \texttt{ConvertAdjacencyListToAdjacencyMatrix} која треба да садржи имплементацију конверзије листе у матрицу суседства.
	\item У Main методи обезбедити начин за тестирање и проверу нове функционалности.
\end{itemize}

\zadatak{2}{Сви карактери првог присутни у другом}
Написати методу која прима два string-а и проверава да ли се сви јединствени карактери из првог string-а појављују барем једном у другом string-у (без обзира на број појављивања).

\textbf{Додатни захтеви за решавање задатка:}
\begin{itemize}
	\item Користити Hash табелу из шаблона за чување броја појављивања сваког карактера.
	\item Имплементацију писати у методи:
	\begin{Verbatim}[fontsize=\small]
public static bool areAllCharsPresent(string str1, string str2)
	\end{Verbatim}
	\item Игнорисати разлику између малих и великих слова.
	\item Уколико једна реч има више истих слова, довољно је постојање само једног карактера у другој речи. (Погледај пример 1)
\end{itemize}

\textbf{Пример:}
\begin{itemize}
	\item \textbf{Улаз 1:} \texttt{str1 = "Aba", str2 = "Balkon"}
	\item \textbf{Излаз 1:} \texttt{true} (Карактери \texttt{"a"} и \texttt{"b"} постоје у другом string-у)
	\item \textbf{Улаз 2:} \texttt{str1 = "Kod", str2 = "Algoritam"}
	\item \textbf{Излаз 2:} \texttt{false} (Карактери \texttt{"d"} и \texttt{"k"} не постоје у другом string-у)
\end{itemize}

\zadatak{3}{Провера "Путања раста" (Root-to-Leaf Max Path)}
У оквиру класе \texttt{BinarySearchTree} написати методу која проверава да ли постоји путања од корена до листа на којој је сваки чвор строго већи од збира свих својих предака на тој путањи.

\textbf{Додатни захтеви за решавање задатка:}
\begin{itemize}
	\item Јавна метода:
	\begin{Verbatim}[fontsize=\small]
public bool existsSuperIncreasingPath()
	\end{Verbatim}
	\item Приватна метода:
	\begin{Verbatim}[fontsize=\small]
private bool checkSuperIncreasingPath(Node current, int currentSum)
	\end{Verbatim}
	\item Подразумева се да стабло има барем два чвора.
	\item Алгоритам ОБАВЕЗНО писати у оквиру класе \texttt{BinarySearchTree}.
\end{itemize}

\textbf{Примери:}

\begin{center}
\begin{minipage}{\textwidth}
\centering
\begin{forest}
for tree={
	circle,
	draw,
	thick,
	fill=gray!20,
	minimum size=8mm,
	inner sep=1pt,
	s sep=8mm,
	l sep=12mm,
	edge={-, thick},
	font=\small
}
[8
	[4
		[2
			[, phantom]
			[3]
		]
		[5]
	]
	[9
		[, phantom]
		[11]
	]
]
\end{forest}

\vspace{0.3cm}
\textbf{Слика 1.} Пример BST-a.
\end{minipage}
\end{center}

\begin{itemize}
	\item \textbf{Улаз 1:} BST са слике 1
	\item \textbf{Излаз 1:} \texttt{false} (11 није веће од 9 + 8)
\end{itemize}

\begin{center}
\begin{minipage}{\textwidth}
\centering
\begin{forest}
for tree={
	circle,
	draw,
	thick,
	fill=gray!20,
	minimum size=8mm,
	inner sep=1pt,
	s sep=8mm,
	l sep=12mm,
	edge={-, thick},
	font=\small
}
[10
	[4
		[2
			[3]
			[, phantom]
		]
		[6
			[5]
			[9]
		]
	]
	[20
		[12
			[, phantom]
			[14]
		]
		[29
			[27]
			[, phantom]
		]
	]
]
\end{forest}

\vspace{0.3cm}
\textbf{Слика 2.} Пример BST-a.
\end{minipage}
\end{center}

\begin{itemize}
	\item \textbf{Улаз 2:} BST са слике 2
	\item \textbf{Излаз 2:} \texttt{false} (29 није веће од 10 + 20)
\end{itemize}

\sablon{Шаблон другог поправног теста фебруар 2026. група А}{2025/Februar/GrupaA/AISP2025_PopravniTest2_Test2_GrupaA.linq}
